\label{chap:req}

In this chapter we outline what the detector was required to do and what restricted us in our
design. Then we consider how it affects society, safety, the environment and the law, along with
the ethical issues it raised, the standards we adhered to, and the planning, risk and cost
management of the project.

\section{Final Specifications and Requirements}
\label{sec:specs}

\subsection{Objectives and specification}

We developed a lie detector for Llama-3.3-70B-Instruct. It reads a response and assigns it a
score, and the higher the score, the more likely the response is to be a lie. In addition to the
detector, we constructed the evaluation to determine how effective it is. Three conditions are
imposed on it by the objectives in \cref{sec:objective}. No labels from the evaluation datasets could be
used to construct it. It had to be tested the same way Liars' Bench tests its own detectors,
otherwise our numbers could not be compared to theirs. And anyone else should be able to
validate our claims from the result files we provide.

\subsection{Constraints}

\begin{description}[leftmargin=1.2em, labelindent=0em, itemsep=2pt]
\item[C1: Model access.] In 16-bit precision the weights of the 70B model occupy about 140\,GB,
  which is far more than we have locally, so we had to run it remotely with NDIF
  \cite{fiottokaufman2024nnsight}, within the limits of their service. The three models of the
  replication (\cref{sec:xmodel}) were run the same way.
\item[C2: Bandwidth.] Download speeds were around 2\,MB/s. At this speed a full per-token cache of
  a single layer, about 60\,GB, would take approximately 8 hours, so we pooled the activations on
  the server whenever possible.
\item[C3: No training of the model.] The model is used exactly as released, with no fine-tuning and
  no gradients.
\item[C4: Licences.] The model is made available by Meta under its Llama 3.3 community licence
  \cite{meta2024llama33}, and the three models of the replication under the Qwen licence agreement
  \cite{qwen2024qwen25}, the Apache 2.0 licence \cite{mistral2025small31} and the Gemma terms of use
  \cite{gemma2025gemma3}. Each dataset is used under the conditions published by its authors
  (\cref{sec:ethics}).
\item[C5: Time.] The final-phase experiments, the analysis and the writing all had to be done within
  one semester (\cref{sec:pm}).
\end{description}

\subsection{Data, software and hardware requirements}

\Cref{tab:software} lists the software and hardware the solution depends on. All software is
open source. The data (five Liars' Bench datasets, RepE, Alpaca, and the persona rollouts) are
described in \cref{sec:data}.

\begin{table}[t]
\centering
\small
\caption{Software and hardware used, with the versions of the analysis environment.}
\label{tab:software}
\begin{tabularx}{\linewidth}{@{}l l X@{}}
\toprule
Component & Version & Role \\
\midrule
Python & 3.12.10 & analysis language \\
PyTorch & 2.7.1 (CPU) & tensors, activation handling \\
NNsight & 0.7.0 & remote tracing of the 70B model on NDIF \\
Transformers & 4.55.2 & tokenisation and chat templates \\
PEFT & 0.15.2 & applying a LoRA adapter inside a remote trace (Gender-Secret) \\
scikit-learn & 1.6.1 & logistic-regression probes, PCA, AUROC \\
NumPy, pandas & 1.26.4, 2.2.2 & numerical work and result tables \\
Matplotlib & 3.9.4 & figures \\
pytest & 8.2.2 & automated tests \\
Tectonic & 0.17.0 & building this report and the paper \\
Git, GitHub & 2.47.1 & version control \\
\midrule
NDIF & n/a & remote access to Llama-3.3-70B-Instruct and Llama-3.1-70B-Instruct, and to the
three models of the replication \\
Local machine & 16\,GB RAM, CPU only & all analysis \\
\bottomrule
\end{tabularx}
\end{table}

\section{Societal Impact}
\label{sec:societal}

\subsection{Society}

Language models can now do things for people. They answer health and legal questions, write
code and use tools. All of this is of little value if we cannot believe what the model is telling us,
either about the world or about itself. A detector that notices when a model says something it
believes is not true helps with this. It can be used to monitor deployed systems, with flagged
responses passed to a human. Before release, it can audit how often a model lies under
pressure. Researchers developing more honest models can also use it to measure their
progress. Since scoring costs almost nothing (\cref{sec:environment}), every response can be checked, even by
organisations that could not afford a second model call for each one.

The detector comes with risks too, and we want to state them openly. The most serious are
false negatives. Our detector misses a lot of lies: averaged over Liars' Bench, even the best
version of our detector catches fewer than half of them at the 1\% threshold (\cref{sec:cmp-benchmark}). It must
never be the sole protection. False positives are also a problem, since every honest answer that
gets flagged costs some trust in the system and the monitor. The 1\% budget holds on benign
chat, but \cref{sec:stats-calibration} shows it does not hold on all datasets. There is also the danger of relying
too heavily on one score to make decisions which people should make. Finally, when a detector
is published it becomes something to outwit, and a model trained against it can learn to lie in
ways the detector does not catch.

\subsection{Health and safety}

The proposed detector offers the greatest benefit to the dataset with the most direct safety
impact, namely \hpkr, which is the dataset that is being studied. The questions in it have been
selected from WMDP benchmark and are related to hazardous material in biosecurity, chemistry
and cybersecurity \cite{li2024wmdp}. Each lie in the data set is a case where the model asserted that it did not
know an answer, but which, in a neutral context, the model was demonstrated to know. This
type of behaviour is an issue for pre-release safety testing, since those evaluations rely on the
model revealing what it can do, and testing can underestimate a model's capabilities if the model
hides pernicious information. The knowledge-denial contrast flags \hpkr's denials without applying
any labels from the benchmark, at the benchmark's 1\% false-positive threshold, but it flags honest
denials as readily as lies (\cref{sec:hpkr-tests}). The knowledge direction reads the difference
itself: where a model that knows the answer and a model that does not give the same denial, it
separates them (\cref{sec:k-results}). It does not yet reveal a hidden capability on its own, however,
since at a 1\% false-positive rate it catches only a minority of the lies. Its reliability should not
be taken for granted either, since \hpkr{} is not a blind test of it, it has been evaluated in full on
only one model, and on the one other model with a Gender-Secret fine-tune it did not detect that
fine-tune's lies (\cref{sec:xm-k,sec:limitations}); it is therefore best used alongside existing
capability evaluations.

\subsection{Legal}

AI regulation is rapidly evolving, and AI providers today are increasingly asked to keep a close
eye on the actions of their systems and to be transparent about them. The European Union's
Artificial Intelligence Act, for instance, has transparency and human-oversight requirements for
some systems \cite{euaiact2024}. One way a provider may fulfil such obligations may include the use of
deception detectors. If they ever are used as evidence of compliance, however, their error rates
must be very carefully characterised, and \cref{sec:stats-calibration} demonstrates that the threshold used
here does not work equally well on all of the data sets. The legal issues that arose through the
project were restricted to licensing, and the model, datasets and code were utilised within the
limits of the licence terms described in \cref{sec:specs,sec:ethics}.

\subsection{Culture}

This study was carried out in English only. This is also true of the datasets and the role-play
prompts and, although less obviously so, of the personas, as the spin doctor, the whistle blower
and the ombudsman are terms of a professional and largely western vocabulary of lies and
honesty. In different languages and cultures, lies and admission/denial of knowledge are
conveyed in different ways, and what is polite in one context may be interpreted as evasive in
another. Such differences cannot be assumed to be captured by a direction coming from English
role play. The detection coverage depends on the type of lies the roles enact, so if one is to be
extended to another language, the personas would need to be created specifically for that
language and culture and a translation of the existing ones would likely not be enough.

\section{Environmental Impact}
\label{sec:environment}

The environmental cost of this work lies almost entirely in running the 70B models remotely; the
analysis runs on a laptop CPU. We estimate it from what was measured, stating every assumption,
because an order-of-magnitude figure is more useful than none \cite{strubell2019energy,patterson2021carbon}.

\subsection{Compute estimate}

A forward pass of a model with $N$ parameters costs about $2N$ floating-point operations per token
\cite{kaplan2020scaling}, about $1.4 \times 10^{11}$ for a 70B model. The largest single measured
activity, the \V{depth-n}-layer sweep, pooled \V{th-sweep-tokens} million tokens of evaluation data
through the model in one pass. Summing the activities (generating and judging the
\V{bank-generated} persona responses, extracting their activations, three full per-token extractions
of the evaluation data, the sweep, and the smaller Gender-Secret, AISI and follow-up runs) gives an
estimate of 20 to 40 million tokens through 70B models over the project. The reads of the final week,
for the knowledge direction and through the fine-tunes, used short prompts or a few hundred
transcripts each and fall within this estimate. The replication of \cref{sec:xmodel} repeated this
work, on fewer datasets, for three more models, two of them less than half the size; it therefore adds
at most three times the estimate, which does not change the comparisons that close this subsection.
\authornote{Replace with the exact total if NDIF usage records are available.}

\begin{table}[t]
\centering
\small
\caption{Order-of-magnitude energy estimate. Every row is an assumption except the first, whose sweep
component is measured and whose total is estimated; the conclusions below hold across the whole range.}
\label{tab:energy}
\begin{tabular}{@{}l l@{}}
\toprule
Quantity & Value \\
\midrule
Tokens through 70B models (sweep measured; total estimated) & $2$--$4 \times 10^{7}$ \\
Operations per token, $2N$ \cite{kaplan2020scaling} & $1.4 \times 10^{11}$ \\
Total operations & $3$--$6 \times 10^{18}$ \\
Effective throughput per accelerator (assumed, $\sim$30\% of peak) & $3 \times 10^{14}$ per second \\
Accelerator-hours & 3--5 \\
Power per accelerator, with data-centre overhead (assumed) & $\sim$0.85\,kW \\
Remote energy, allowing up to $10\times$ serving inefficiency & $\sim$2--45\,kWh \\
Local analysis (laptop, $\sim$100\,h at $\sim$40\,W, assumed) & $\sim$4\,kWh \\
Emissions at 0.4\,kg\,CO$_2$e per kWh (assumed grid intensity) & $\sim$2.5--20\,kg CO$_2$e \\
\bottomrule
\end{tabular}
\end{table}

\Cref{tab:energy} gives the resulting range: a few to a few tens of kilowatt-hours, and at most
twenty kilograms of CO$_2$-equivalent, at most about as much as driving a petrol car a hundred
kilometres. Two comparisons put this in context. Training a model of this size takes millions of
accelerator-hours \cite{grattafiori2024llama}; this study used a vanishing fraction of that, and
reused a released model rather than training one. And the persona bank is built once and reused for
every dataset, whereas a supervised approach would repeat labelling and training for each new
domain.

\subsection{Design decisions taken to minimise its impact}

The following engineering decisions reduced the compute required, and are part of the design:
\begin{itemize}[leftmargin=1.2em, itemsep=2pt]
\item \emph{One pass for all layers.} The sweep loaded all \V{th-sweep-layers-new} new layers (those beyond the
  three first extracted) in one pass for each of the examples, because the remote pass calculated all the layers anyway. The cost
  of reading them one by one would have been \V{th-sweep-layers-new} times as much.
\item \emph{Pooling on the server.} The means over tokens have been computed remotely, hence a layer
  costs $\sim$300\,MB to transfer compared to $\sim$60\,GB.
\item \emph{Resumable, cached extraction.} Completed requests were saved and it was never necessary
  to repeat completed work when the remote 70B model was actually interrupted, which was often the
  case.
\item \emph{CPU-only analysis.} Once the extraction is done, all experiments are performed on a
  laptop and all \V{depth-n} layers can be analysed within an hour on 4 cores.
\end{itemize}

\subsection{Sustainability of the solution in use}

The detector is environmentally friendly. It takes one dot product of two 8192-dimensional
(about $1.6 \times 10^{4}$ operations) vectors to score a response, using activations the model has
already computed as it generates the response. A black box judge, however, has to run a
second time to read a 500 token response, which for a 70B judge is about $7 \times 10^{13}$ operations,
over a billion times as many. With an assistant that is under constant scrutiny, the difference will
be between a monitor that adds almost no energy and one that will roughly double the serving
cost when the judge is the same size as the assistant.

\section{Ethical Issues}
\label{sec:ethics}

\emph{Hazardous content.} The development of \hpkr{} is based on hazardous-knowledge questions
\cite{li2024wmdp}. Only the existing responses in the benchmark were used and the model's activations were
read on the existing responses. No hazardous information was created, extended or shared,
and there is no information in this report or in the released files beyond that which is already
provided to the public in the benchmark. Our own role-play questions are about disclosure in the
day-to-day professional setting, and do not include anything dangerous.

\emph{Dual use.} Lying detection research could also be employed to help a model evade detection.
We felt it was more important to have a public and correct record of what these detectors can
and cannot do. The outline of the detector design is already available \cite{mahadik2026persona} and
knowledge that the detector has weaknesses, such as when both classes deny knowledge
(\cref{sec:hpkr-tests}), helps the defender more than the evader. No method designed for evasion is
released.

\emph{Honest reporting.} The project was dedicated to accurate reporting of the results, good and bad.
Claims about all layers were checked against the \V{depth-n}-layer sweep, and those that held only at
the three layers first extracted are stated in their narrower form rather than dropped silently
(\cref{sec:analysis,sec:adjustments}). A
collaborator's activation-extraction pipeline gave a set of intermediate results that were found to
come from an error in the pipeline. Those results have been discarded and are not reported or relied
on in this thesis. Most numbers are read directly from result files, and the few that are not are listed
with their sources in \cref{sec:code}.

\emph{Use conditions and licences.} Each model/dataset was used under the terms of its publication. The AISI rollouts \cite{aisi_lie_rollouts} are published under a licence
marked ``other'', without specifying the terms of the licence. We used it only for research
purposes, we did not redistribute it and we used one fixed revision.
\authornote{Confirm the AISI dataset's terms with its publisher before publishing results derived
from it.}

\emph{Attribution.} The persona bank has been created by one member of the team and the analysis in
this thesis by the authors (\cref{sec:pm}).
\authornote{Name the member who built the persona bank here, and fill in each member's role in
the Resources subsection of the Project Management Plan.} The use of other methods and data is indicated by
citation.


\section{Project Management Plan}
\label{sec:pm}

\subsection{Phases and schedule}

The thesis ran over three phases, following the course structure. \Cref{tab:schedule} gives the
schedule; the two pre-thesis phases are documented in their own reports.

\begin{table}[t]
\centering
\small
\caption{Project schedule.}
\label{tab:schedule}
\begin{tabularx}{\linewidth}{@{}l l X@{}}
\toprule
Phase & Period & Work \\
\midrule
Pre-Thesis 1 & \fillin{semester} & Problem formulation, literature review, background study \\
Pre-Thesis 2 & \fillin{semester} & Preliminary design, choice of benchmark and model, access to NDIF \\
\midrule
Final & July 2026 & Repository set-up; the benchmark's probe pipeline vendored and adapted \\
      & July--Aug.\ 2026 & Persona bank construction: prompts, generation, judging, vectors \\
      & 4--5 Aug.\ & Extraction infrastructure: resumable cache, token-budget batching, NDIF runner \\
      & 10 Aug.\ & Layer 40, the model's middle layer, enters the pipeline \\
      & 15 Aug.\ & Probe code reads layer 40 by default \\
      & 16 Aug.\ & Layer sweep \\
      & 18 Aug.\ & Persona geometry; transfer experiment; sycophancy control \\
      & 20 Aug.\ & Evaluation moves to layer 40; knowledge-denial personas; zero-shot evaluation \\
      & 21 Aug.\ & Layer plan fixed (\cref{sec:layer-choice}); layer 44; Follow-Up Probe reproduced;
                    $k=3$ subspace \\
      & 22 Aug.\ & Token reads; Gender-Secret; AISI baseline; paper draft \\
      & 23 Aug.\ & \V{depth-n}-layer sweep; unified-axis control \\
      & 24 Aug.\ & Paper revision; this report \\
      & 25--26 Aug.\ & Wording-controlled \hpkr{} slices; forced-denial reads: the wording and
                    chemistry knowledge directions \\
      & 28 Aug.\ & The knowledge direction; Gender-Secret and Soft-Trigger read through their
                    fine-tunes \\
      & 28--29 Aug.\ & Robustness, replication and operating points of the knowledge direction \\
      & 6--12 Oct.\ & Persona banks for Qwen2.5-72B, Mistral-Small-3.1-24B and Gemma-3-27B; the
                    benchmark's probes reproduced on each \\
      & 13--15 Oct.\ & Every study but steering repeated on the three models; paired intervals; the
                    template-free K (\cref{sec:xmodel}) \\
\bottomrule
\end{tabularx}
\end{table}

\subsection{Resources}

\begin{itemize}[leftmargin=1.2em, itemsep=2pt]
\item \emph{People.} Md.\ Abdullah Al Zami, Suraiya Shanaz and Mahdi Hasan.
  \authornote{Add each member's role, e.g.\ persona bank construction; extraction and analysis;
  writing.}
\item \emph{Compute.} Remote access to Llama-3.3-70B-Instruct and Llama-3.1-70B-Instruct, and for
  the replication to Qwen2.5-72B, Mistral-Small-3.1-24B and Gemma-3-27B, through NDIF; one 16\,GB
  laptop for all analysis.
\item \emph{Data.} The public Liars' Bench, RepE and Alpaca datasets; the AISI rollouts; the persona
  rollouts generated by the team.
\item \emph{Software.} The open-source stack of \cref{tab:software}.
\end{itemize}

\subsection{Monitoring}

Progress was tracked through the repository. Each experiment's results were written to a folder
with a summary of what it showed. When a result contradicted an earlier summary, the earlier summary
was corrected in the same step, so the written record never lagged the evidence.

\section{Risk Management}
\label{sec:risk}

\Cref{tab:risk} is the project's risk register. Every risk in it occurred at least once; the table
records what happened and how it was handled, which is more informative than a list of risks
anticipated in the abstract.

\begin{table}[t]
\centering
\small
\caption{Risk register: risks met during the project, with their mitigation and outcome.}
\label{tab:risk}
\begin{tabularx}{\linewidth}{@{}>{\raggedright\arraybackslash}p{3.0cm} c >{\raggedright\arraybackslash}X >{\raggedright\arraybackslash}p{3.1cm}@{}}
\toprule
Risk & Impact & Mitigation & Outcome \\
\midrule
Remote out-of-memory on long examples & high & Token-budget batching that halves the batch after
repeated failures; rows too long for any batch dropped and counted & All datasets extracted \\
Local memory ($16$\,GB) too small for large probes & high & A streaming logistic regression that
solves scikit-learn's exact problem block by block, tested for equivalence & AISI probe trained \\
Slow download ($\sim$2\,MB/s) & high & Pooling on the server; all layers in one pass & A layer in
$\sim$300\,MB instead of $\sim$60\,GB \\
Interrupted remote jobs & medium & Every completed request saved; jobs resume from the last one &
No work repeated \\
Erroneous intermediate results from a collaborator's pipeline & high & Results re-derived in an
independently tested pipeline; the erroneous set discarded & None of it used \\
Claims that do not generalise beyond the layers studied & high & A \V{depth-n}-layer sweep, and
build-time assertions on every ``at every layer'' claim & Several claims narrowed to the layers where
they hold \\
Overwriting hand-written results & medium & The analysis script refuses to overwrite a file it did
not write & Three files restored \\
Unclear dataset terms & medium & Research use only, no redistribution, revision pinned & Flagged
for confirmation \\
Remote out-of-memory on Gemma's long rows and large vocabulary & high & Requests capped by tokens,
down to one row per request, and shorter requests wherever logits are computed & Every study completed \\
A model that does not play the roles (Qwen) & high & Its lenient bank used, since the strict one lacks
11 roles; reported as a finding, prompts not tuned & Persona claims fail on Qwen \\
\bottomrule
\end{tabularx}
\end{table}

\section{Economic Analysis}
\label{sec:econ}

\subsection{Costs}

The project was virtually free, especially in terms of direct monetary costs. We used a laptop
computer that we already had, and all the software and datasets used are free. Additionally,
accessing the model remotely through NDIF cost us nothing. The actual cost was the time that
the team devoted. \authornote{Estimate the team's working hours and, if required, value them at an
appropriate rate.}

It is still advisable to estimate how much the remote compute would cost if it had to be
purchased commercially, thereby requiring renting of accelerators large enough for a 70B
model. \Cref{tab:energy} estimates the work at 3--5 accelerator-hours at full efficiency, and up to 10
times this when allowing for serving overhead, for a total range of 3--50 hours. If accelerators
cost \$2 to \$4 per accelerator-hour, this will cost between about \$6 and \$200. On top of this cost
there would be other costs involved, including the engineering process, minimum rental periods
and storage, as a full activation cache takes around 60\,GB per layer.
\authornote{Update with current market prices if required.}

\subsection{Cost--benefit of the solution}

A comparison of the three families of detectors is given in \cref{tab:costbenefit}. The economic case
for a label-free probe is based primarily on two expenses which it does not incur.

Labelling is the first of these. A supervised probe requires labelled lies for each new domain to
which it may be applied, and labelled lies are not readily available as whether a statement is a
lie depends on what the model believes, which cannot be determined by whether the statement
is true. For this reason Liars' Bench runs a separate check of the model's knowledge for each
label. The persona detector's bank has to be constructed once, in our case from \V{bank-generated}
role-play generations and their judgements, and then it requires no labels for any new set of data.

Inference costs are the other consideration. A probe adds almost nothing to the cost of serving,
while a black-box judge needs a second large-model call for each response and could double that cost
(\cref{sec:environment}).

\begin{table}[t]
\centering
\small
\caption{Qualitative cost--benefit comparison of the three detector families.}
\label{tab:costbenefit}
\begin{tabularx}{\linewidth}{@{}>{\raggedright\arraybackslash}p{3.3cm} >{\raggedright\arraybackslash}X >{\raggedright\arraybackslash}X >{\raggedright\arraybackslash}X@{}}
\toprule
 & Persona contrast (this work) & Supervised probe & Black-box judge \\
\midrule
Set-up cost & one persona bank, reused & labelled lies per domain & prompt design \\
Labels per new domain & none & many & none \\
Cost per response & one dot product & one dot product & one large-model call \\
Needs model internals & yes & yes & no \\
Transfer across datasets & highest of the designs compared at layer~40, with row-level intervals
only (\cref{sec:transfer}) & weak & not evaluated here \\
\bottomrule
\end{tabularx}
\end{table}

The major constraint is access. A white-box detector requires access to the model's internals,
which is available to anyone who is deploying an open-weight model and not to anyone who is
using a closed API. When the internals can be accessed, the lowest cost method per response
and per new domain is the persona contrast, and in \cref{chap:results} it also has the highest transfer across datasets
of the designs compared, although that margin carries row-level intervals only and a probe trained
on several pooled datasets nearly matches it once \hpkr{} is left out
(\cref{sec:stats-nulls,sec:transfer}).
